\begin{section}{Metric Spaces}

Let $X$ be a set. A function $d: X \times X \to \R$ is called a \textbf{metric} on $X$ if for any $p, q \in X$, we have
\begin{enumerate}[label=\textbf{(\Alph*)}]
    \item $d(p, q) \geq 0$ and $d(p, q) = 0 \iff p = q$.
    \item $d(p, q) = d(q, p)$.
    \item $d(p, r) \leq d(p, q) + d(q, r)$ for any $r \in X$.
\end{enumerate}
\textbf{(C)} is called the \textbf{triangle inequality}. The pair $(X, d)$ is called a \textbf{metric space}.

\bigskip

\begin{theorem} \label{thm:StandardMetricOnRn}
    Let $d: \R^{n} \times \R^{n} \to \R$ be a function defined by
    \[
    d(\mathbf{x}, \mathbf{y}) = \norm{\mathbf{x} - \mathbf{y}} = \sqrt{\sum_{k = 1}^{n} (x_{k} - y_{k})^{2}}
    \]
    for all $\mathbf{x}, \mathbf{y} \in \R^{n}$. Then $d$ is a metric on $\R^{n}$.
    Especially, $d$ is called the \textbf{standard metric} on $\R^{n}$.
\end{theorem}

\begin{proof}
    Let $\mathbf{x}, \mathbf{y}, \mathbf{z} \in \R^{n}$ be given.
    \begin{itemize}
        \item $d(\mathbf{x}, \mathbf{y}) \geq 0$ is obvious. Also, we have
        \begin{align*}
            d(\mathbf{x}, \mathbf{y}) = 0 &\iff \norm{\mathbf{x} - \mathbf{y}} = 0 \\
            &\iff x_{k} = y_{k} \text{ for all } k = 1, 2, \dots, n \\
            &\iff \mathbf{x} = \mathbf{y}.
        \end{align*}
        \item $d(\mathbf{x}, \mathbf{y}) = d(\mathbf{y}, \mathbf{x})$ is obvious.
        \item By \Cref{thm:PropertiesOfInnerProductNormDistance}, we have
        \begin{align*}
            d(\mathbf{x}, \mathbf{z}) &= \norm{\mathbf{x} - \mathbf{z}} =
            \norm{(\mathbf{x} - \mathbf{y}) + (\mathbf{y} - \mathbf{z})} \\
            &\leq \norm{\mathbf{x} - \mathbf{y}} + \norm{\mathbf{y} - \mathbf{z}}
            = d(\mathbf{x},\mathbf{y}) + d(\mathbf{y},\mathbf{z})
        \end{align*}
    \end{itemize}
    Therefore, $d$ is a metric on $\R^{n}$.
\end{proof}

\bigskip

\begin{definition} \label{def:NeighborhoodOpenBall}
    Let $(X, d)$ be a metric space. For any $p \in X$ and $r > 0$,
    we define the \textbf{neighborhood} of $p$ with radius $r$ as
    \[
    N_{r}(p) \coloneqq \setbuilder{q \in X}{d(p, q) < r}.
    \]
    Especially, if $X = \R^{n}$ and $d$ is the standard metric on $\R^{n}$, then $N_{r}(p)$ is called the
    \textbf{open ball} of radius $r$ centered at $p$, which is denoted by $B_{r}(p)$.
\end{definition}

\bigskip

\begin{definition} \label{def:LimitIsolatedInteriorPoints}
    Let $(X, d)$ be a metric space and $E \subseteq X$. Let $p \in X$ be given.
    \begin{itemize}
        \item $p$ is called an \textbf{limit point} of $E$
        if for any $r > 0$, we have $E \cap (N_{r}(p) \setminus \set{p}) \neq \emptyset$.
        \item $p$ is called a \textbf{isolated point} of $E$
        if there exists some $r > 0$ such that $N_{r}(p) \cap E = \set{p}$.
        \item $p$ is called an \textbf{interior point} of $E$
        if there exists some $r > 0$ such that $N_{r}(p) \subseteq E$.
    \end{itemize}
\end{definition}

\bigskip

\begin{definition} \label{def:OpenClosedSets}
    Let $(X, d)$ be a metric space and $E \subseteq X$.
    \begin{itemize}
        \item $E$ is called \textbf{open} if any point of $E$ is an interior point of $E$.
        \item $E$ is called \textbf{closed} if any limit point of $E$ belongs to $E$.
    \end{itemize}
\end{definition}

\bigskip

\begin{theorem} \label{thm:NeighborhoodIsOpen}
    Let $(X, d)$ be a metric space. For any $p \in X$ and $r > 0$, the neighborhood $N_{r}(p)$ is open.
\end{theorem}

\begin{proof}
    Let $q \in N_{r}(p)$ be given. Then we have $d(p, q) < r$. Put $\veps = r - d(p, q) > 0$.
    Let $s \in N_{\veps}(q)$ be given. Then we have $d(q, s) < \veps$ and
    \[
    d(p, s) \leq d(p, q) + d(q, s) < d(p, q) + \veps = r.
    \]
    Thus $s \in N_{r}(p)$, which implies that $N_{\veps}(q) \subseteq N_{r}(p)$. Since $q$ is arbtrary, $N_{r}(p)$ is open.
\end{proof}

\end{section}